\documentclass{amsart}
% For dealing with references we use the comment environment
\usepackage{verbatim}
\newenvironment{reference}{\comment}{\endcomment}
%\newenvironment{reference}{}{}
\newenvironment{slogan}{\comment}{\endcomment}
\newenvironment{history}{\comment}{\endcomment}

% For commutative diagrams we use Xy-pic
\usepackage[all]{xy}

% We use 2cell for 2-commutative diagrams.
\xyoption{2cell}
\UseAllTwocells

% We use multicol for the list of chapters between chapters
\usepackage{multicol}

% This is generally recommended for better output
\usepackage{lmodern}
\usepackage[T1]{fontenc}

% For cross-file-references

% Package for hypertext links:
\usepackage{hyperref}

% For any local file, say "hello.tex" you want to link to please

% Theorem environments.
%
\theoremstyle{plain}
\newtheorem{theorem}[subsection]{Theorem}
\newtheorem{proposition}[subsection]{Proposition}
\newtheorem{lemma}[subsection]{Lemma}

\theoremstyle{definition}
\newtheorem{definition}[subsection]{Definition}
\newtheorem{example}[subsection]{Example}
\newtheorem{exercise}[subsection]{Exercise}
\newtheorem{situation}[subsection]{Situation}

\theoremstyle{remark}
\newtheorem{remark}[subsection]{Remark}
\newtheorem{remarks}[subsection]{Remarks}

\numberwithin{equation}{subsection}

% Macros
%
\def\lim{\mathop{\mathrm{lim}}\nolimits}
\def\colim{\mathop{\mathrm{colim}}\nolimits}
\def\Spec{\mathop{\mathrm{Spec}}}
\def\Hom{\mathop{\mathrm{Hom}}\nolimits}
\def\Ext{\mathop{\mathrm{Ext}}\nolimits}
\def\SheafHom{\mathop{\mathcal{H}\!\mathit{om}}\nolimits}
\def\SheafExt{\mathop{\mathcal{E}\!\mathit{xt}}\nolimits}
\def\Sch{\mathit{Sch}}
\def\Mor{\mathop{\mathrm{Mor}}\nolimits}
\def\Ob{\mathop{\mathrm{Ob}}\nolimits}
\def\Sh{\mathop{\mathit{Sh}}\nolimits}
\def\NL{\mathop{N\!L}\nolimits}
\def\CH{\mathop{\mathrm{CH}}\nolimits}
\def\proetale{{pro\text{-}\acute{e}tale}}
\def\etale{{\acute{e}tale}}
\def\QCoh{\mathit{QCoh}}
\def\Ker{\mathop{\mathrm{Ker}}}
\def\Im{\mathop{\mathrm{Im}}}
\def\Coker{\mathop{\mathrm{Coker}}}
\def\Coim{\mathop{\mathrm{Coim}}}

% Boxtimes
%
\DeclareMathSymbol{\boxtimes}{\mathbin}{AMSa}{"02}

%
% Macros for moduli stacks/spaces
%
\def\QCohstack{\mathcal{QC}\!\mathit{oh}}
\def\Cohstack{\mathcal{C}\!\mathit{oh}}
\def\Spacesstack{\mathcal{S}\!\mathit{paces}}
\def\Quotfunctor{\mathrm{Quot}}
\def\Hilbfunctor{\mathrm{Hilb}}
\def\Curvesstack{\mathcal{C}\!\mathit{urves}}
\def\Polarizedstack{\mathcal{P}\!\mathit{olarized}}
\def\Complexesstack{\mathcal{C}\!\mathit{omplexes}}
% \Pic is the operator that assigns to X its picard group, usage \Pic(X)
% \Picardstack_{X/B} denotes the Picard stack of X over B
% \Picardfunctor_{X/B} denotes the Picard functor of X over B
\def\Pic{\mathop{\mathrm{Pic}}\nolimits}
\def\Picardstack{\mathcal{P}\!\mathit{ic}}
\def\Picardfunctor{\mathrm{Pic}}
\def\Deformationcategory{\mathcal{D}\!\mathit{ef}}

\setlength{\emergencystretch}{3em}
\begin{document}
\title[Stacks Project, Tag 0728]{2869 --- A surjective separated etale map of algebraic spaces yields an integral alternating Cech spectral sequence on unordered distinct-point configurations with sign local coefficients}
\maketitle
\noindent Copyright (C) 2005--2025 Johan de Jong. Stacks Project authors. Source: \href{https://stacks.math.columbia.edu/tag/0728}{Tag 0728}.

\noindent Editorial difficulty note (not author text). Difficulty H2: The construction replaces ordered etale Cech intersections by unordered configurations while retaining integral coefficients and the sign local system. Its proof must identify exterior-power sheaves with configuration-space sign coinvariants, build the character line sheaf on a finite-group quotient, and compare two injective double-complex spectral sequences. This geometric and derived construction goes beyond a routine Cech-cover computation, and does not assume exactness of group anti-invariants..

\noindent Editorial provenance note (not author text). History: excerpt from revision \texttt{a0444\allowbreak 6e57e\allowbreak c1fbc\allowbreak 252a8\allowbreak 71afc\allowbreak ec775\allowbreak 2fb28\allowbreak 07b14}, spaces-cohomology.tex, lines 999--1072. Reference presentation was adapted; original-author context and core support blocks were retained or added. The mathematical wording is unchanged. Licensed under GNU FDL 1.2 or later; no invariant sections or cover texts. See ../licenses/COPYING. Context blocks reproduce author definitions and supporting results; they are not additional counted problems.

\begin{definition}
\label{definition-alternating-cech-complex}
Let $S$ be a scheme. Let $f : U \to X$ be a surjective \'etale morphism
of algebraic spaces over $S$. Let $\mathcal{F}$ be an object of
$\textit{Ab}(X_\etale)$. The
{\it alternating {\v C}ech complex}\footnote{This may be nonstandard notation}
$\check{\mathcal{C}}^\bullet_{alt}(f, \mathcal{F})$
associated to $\mathcal{F}$ and $f$ is the complex
$$
\Hom(K^0, \mathcal{F}) \to \Hom(K^1, \mathcal{F}) \to
\Hom(K^2, \mathcal{F}) \to \ldots
$$
with Hom groups computed in $\textit{Ab}(X_\etale)$.
\end{definition}

\section{The alternating {\v C}ech complex}
\label{section-alternating-cech}

\noindent
Let $S$ be a scheme. Let $f : U \to X$ be an \'etale morphism of algebraic
spaces over $S$. The functor
$$
j : U_{spaces, \etale} \longrightarrow X_{spaces, \etale},\quad
V/U \longmapsto V/X
$$
induces an equivalence of $U_{spaces, \etale}$ with the localization
$X_{spaces, \etale}/U$, see
Properties of Spaces, Section \href{https://stacks.math.columbia.edu/tag/04LX}{section localize (Tag 04LX)}.
Hence there exist functors
$$
f_! : \textit{Ab}(U_\etale) \longrightarrow
\textit{Ab}(X_\etale),\quad
f_! : \textit{Mod}(\mathcal{O}_U) \longrightarrow \textit{Mod}(\mathcal{O}_X),
$$
which are left adjoint to
$$
f^{-1} : \textit{Ab}(X_\etale) \longrightarrow
\textit{Ab}(U_\etale),\quad
f^* : \textit{Mod}(\mathcal{O}_X) \longrightarrow \textit{Mod}(\mathcal{O}_U)
$$
see
Modules on Sites, Section \href{https://stacks.math.columbia.edu/tag/03DH}{section localize (Tag 03DH)}.
Warning: This functor, a priori, has
nothing to do with cohomology with compact supports!
We dubbed this functor ``extension by zero'' in the reference above.
Note that the two versions of $f_!$ agree as $f^* = f^{-1}$ for
sheaves of $\mathcal{O}_X$-modules.

\medskip\noindent
As we are going to use this construction below let us recall some of its
properties. Given an abelian sheaf $\mathcal{G}$ on $U_\etale$
the sheaf $f_!$ is the sheafification of the presheaf
$$
V/X \longmapsto
f_!\mathcal{G}(V) =
\bigoplus\nolimits_{\varphi \in \Mor_X(V, U)}
\mathcal{G}(V \xrightarrow{\varphi} U),
$$
see
Modules on Sites, Lemma \href{https://stacks.math.columbia.edu/tag/03DI}{lemma extension by zero (Tag 03DI)}.
Moreover, if $\mathcal{G}$ is an $\mathcal{O}_U$-module, then $f_!\mathcal{G}$
is the sheafification of the exact same presheaf of abelian groups which
is endowed with an $\mathcal{O}_X$-module structure in an obvious way
(see loc.\ cit.). Let $\overline{x} : \Spec(k) \to X$
be a geometric point. Then there is a canonical identification
$$
(f_!\mathcal{G})_{\overline{x}} =
\bigoplus\nolimits_{\overline{u}} \mathcal{G}_{\overline{u}}
$$
where the sum is over all $\overline{u} : \Spec(k) \to U$ such that
$f \circ \overline{u} = \overline{x}$, see
Modules on Sites, Lemma \href{https://stacks.math.columbia.edu/tag/0710}{lemma stalk j shriek (Tag 0710)}
and
Properties of Spaces, Lemma
\href{https://stacks.math.columbia.edu/tag/04K6}{lemma points small etale site (Tag 04K6)}.
In the following we are going to study the sheaf $f_!\underline{\mathbf{Z}}$.
Here $\underline{\mathbf{Z}}$ denotes the constant sheaf on
$X_\etale$ or $U_\etale$.

\begin{lemma}
\label{lemma-product-is-tensor-product}
Let $S$ be a scheme. Let $f_i : U_i \to X$ be \'etale morphisms
of algebraic spaces over $S$. Then there are isomorphisms
$$
f_{1, !}\underline{\mathbf{Z}} \otimes_{\mathbf{Z}}
f_{2, !}\underline{\mathbf{Z}}
\longrightarrow
f_{12, !}\underline{\mathbf{Z}}
$$
where $f_{12} : U_1 \times_X U_2 \to X$ is the structure morphism
and
$$
(f_1 \amalg f_2)_! \underline{\mathbf{Z}}
\longrightarrow
f_{1, !}\underline{\mathbf{Z}} \oplus
f_{2, !}\underline{\mathbf{Z}}
$$
\end{lemma}

\begin{proof}
Once we have defined the map it will be an isomorphism by our description
of stalks above. To define the map it suffices to work on the level of
presheaves. Thus we have to define a map
$$
\left(\bigoplus\nolimits_{\varphi_1 \in \Mor_X(V, U_1)} \mathbf{Z}\right)
\otimes_{\mathbf{Z}}
\left(\bigoplus\nolimits_{\varphi_2 \in \Mor_X(V, U_2)} \mathbf{Z}\right)
\longrightarrow
\bigoplus\nolimits_{\varphi \in \Mor_X(V, U_1 \times_X U_2)}
\mathbf{Z}
$$
We map the element $1_{\varphi_1} \otimes 1_{\varphi_2}$ to the element
$1_{\varphi_1 \times \varphi_2}$ with obvious notation. We omit the proof
of the second equality.
\end{proof}

\noindent
Another important feature is the trace map
$$
\text{Tr}_f : f_!\underline{\mathbf{Z}} \longrightarrow \underline{\mathbf{Z}}.
$$
The trace map is adjoint to the
map $\mathbf{Z} \to f^{-1}\underline{\mathbf{Z}}$ (which is an isomorphism).
If $\overline{x}$ is above, then $\text{Tr}_f$ on stalks at $\overline{x}$
is the map
$$
(\text{Tr}_f)_{\overline{x}} :
(f_!\underline{\mathbf{Z}})_{\overline{x}} =
\bigoplus\nolimits_{\overline{u}} \mathbf{Z}
\longrightarrow
\mathbf{Z} = \underline{\mathbf{Z}}_{\overline{x}}
$$
which sums the given integers. This is true because it is adjoint to the map
$1 : \mathbf{Z} \to f^{-1}\underline{\mathbf{Z}}$. In particular, if
$f$ is surjective as well as \'etale then $\text{Tr}_f$ is surjective.

\medskip\noindent
Assume that $f : U \to X$ is a surjective \'etale
morphism of algebraic spaces. Consider the {\it Koszul complex}
associated to the trace map we discussed above
$$
\ldots \to \wedge^3f_!\underline{\mathbf{Z}} \to
\wedge^2f_!\underline{\mathbf{Z}} \to f_!\underline{\mathbf{Z}} \to
\underline{\mathbf{Z}} \to 0
$$
Here the exterior powers are over the sheaf of rings $\underline{\mathbf{Z}}$.
The maps are defined by the rule
$$
e_1 \wedge \ldots \wedge e_n \longmapsto
\sum\nolimits_{i = 1, \ldots, n} (-1)^{i + 1}
\text{Tr}_f(e_i)
e_1 \wedge \ldots \wedge \widehat{e_i} \wedge \ldots \wedge e_n
$$
where $e_1, \ldots, e_n$ are local sections of $f_!\underline{\mathbf{Z}}$.
Let $\overline{x}$ be a geometric point of $X$ and set
$M_{\overline{x}} = (f_!\underline{\mathbf{Z}})_{\overline{x}} =
\bigoplus_{\overline{u}} \mathbf{Z}$. Then the stalk of the complex above at
$\overline{x}$ is the complex
$$
\ldots \to \wedge^3 M_{\overline{x}} \to \wedge^2 M_{\overline{x}}
\to M_{\overline{x}} \to \mathbf{Z} \to 0
$$
which is exact because $M_{\overline{x}} \to \mathbf{Z}$ is surjective, see
More on Algebra, Lemma \href{https://stacks.math.columbia.edu/tag/0626}{lemma homotopy koszul abstract (Tag 0626)}.
Hence if we let $K^\bullet = K^\bullet(f)$ be the complex with
$K^i = \wedge^{i + 1}f_!\underline{\mathbf{Z}}$, then we obtain a
quasi-isomorphism
\begin{equation}
\label{equation-quasi-isomorphism}
K^\bullet \longrightarrow \underline{\mathbf{Z}}[0]
\end{equation}
We use the complex $K^\bullet$ to define what we call
the alternating {\v C}ech complex associated to $f : U \to X$.



\begin{lemma}
\label{lemma-alternating-cech-to-cohomology}
Let $S$ be a scheme. Let $f : U \to X$ be a surjective \'etale morphism
of algebraic spaces over $S$. Let $\mathcal{F}$ be an object of
$\textit{Ab}(X_\etale)$. There exists a canonical map
$$
\check{\mathcal{C}}^\bullet_{alt}(f, \mathcal{F})
\longrightarrow
R\Gamma(X, \mathcal{F})
$$
in $D(\textit{Ab})$. Moreover, there is a spectral sequence with $E_1$-page
$$
E_1^{p, q} =
\Ext_{\textit{Ab}(X_\etale)}^q(K^p, \mathcal{F})
$$
converging to $H^{p + q}(X, \mathcal{F})$ where
$K^p = \wedge^{p + 1}f_!\underline{\mathbf{Z}}$.
\end{lemma}

\begin{proof}
Recall that we have the quasi-isomorphism
$K^\bullet \to \underline{\mathbf{Z}}[0]$, see
(\ref{equation-quasi-isomorphism}).
Choose an injective resolution $\mathcal{F} \to \mathcal{I}^\bullet$
in $\textit{Ab}(X_\etale)$. Consider the double complex
$\Hom(K^\bullet, \mathcal{I}^\bullet)$ with terms
$\Hom(K^p, \mathcal{I}^q)$. The differential
$d_1^{p, q} : A^{p, q} \to A^{p + 1, q}$
is the one coming from the differential $K^{p + 1} \to K^p$
and the differential $d_2^{p, q} : A^{p, q} \to A^{p, q + 1}$ is the
one coming from the differential
$\mathcal{I}^q \to \mathcal{I}^{q + 1}$.
Denote $\text{Tot}(\Hom(K^\bullet, \mathcal{I}^\bullet))$
the associated total complex, see
Homology, Section \href{https://stacks.math.columbia.edu/tag/0FNB}{section double complexes (Tag 0FNB)}.
We will use the two spectral
sequences $({}'E_r, {}'d_r)$ and $({}''E_r, {}''d_r)$
associated to this double complex, see
Homology, Section \href{https://stacks.math.columbia.edu/tag/012X}{section double complex (Tag 012X)}.

\medskip\noindent
Because $K^\bullet$ is a resolution of $\underline{\mathbf{Z}}$
we see that the complexes
$$
\Hom(K^\bullet, \mathcal{I}^q) :
\Hom(K^0, \mathcal{I}^q) \to
\Hom(K^1, \mathcal{I}^q) \to
\Hom(K^2, \mathcal{I}^q) \to \ldots
$$
are acyclic in positive degrees and have $H^0$ equal to
$\Gamma(X, \mathcal{I}^q)$. Hence by
Homology, Lemma \href{https://stacks.math.columbia.edu/tag/0133}{lemma double complex gives resolution (Tag 0133)}
the natural map
$$
\mathcal{I}^\bullet(X) \longrightarrow
\text{Tot}(\Hom(K^\bullet, \mathcal{I}^\bullet))
$$
is a quasi-isomorphism of complexes of abelian groups. In particular
we conclude that
$H^n(\text{Tot}(\Hom(K^\bullet, \mathcal{I}^\bullet))) = H^n(X, \mathcal{F})$.

\medskip\noindent
The map $\check{\mathcal{C}}^\bullet_{alt}(f, \mathcal{F}) \to
R\Gamma(X, \mathcal{F})$ of the lemma is the composition of
$\check{\mathcal{C}}^\bullet_{alt}(f, \mathcal{F}) \to
\text{Tot}(\Hom(K^\bullet, \mathcal{I}^\bullet))$
with the inverse of the displayed quasi-isomorphism.

\medskip\noindent
Finally, consider the spectral sequence $({}'E_r, {}'d_r)$.
We have
$$
E_1^{p, q} = q\text{th cohomology of }
\Hom(K^p, \mathcal{I}^0) \to
\Hom(K^p, \mathcal{I}^1) \to
\Hom(K^p, \mathcal{I}^2) \to \ldots
$$
This proves the lemma.
\end{proof}

\noindent
It follows from the lemma that it is important to understand the
ext groups $\Ext_{\textit{Ab}(X_\etale)}(K^p, \mathcal{F})$,
i.e., the right derived functors of
$\mathcal{F} \mapsto \Hom(K^p, \mathcal{F})$.

\begin{lemma}
\label{lemma-compute}
Let $S$ be a scheme. Let $f : U \to X$ be a surjective, \'etale, and separated
morphism of algebraic spaces over $S$. For $p \geq 0$ set
$$
W_p = U \times_X \ldots \times_X U \setminus \text{all diagonals}
$$
where the fibre product has $p + 1$ factors.
There is a free action of $S_{p + 1}$ on $W_p$ over $X$ and
$$
\Hom(K^p, \mathcal{F}) = S_{p + 1}\text{-anti-invariant elements of }
\mathcal{F}(W_p)
$$
functorially in $\mathcal{F}$ where
$K^p = \wedge^{p + 1}f_!\underline{\mathbf{Z}}$.
\end{lemma}

\begin{proof}
Because $U \to X$ is separated the diagonal $U \to U \times_X U$ is a
closed immersion. Since $U \to X$ is \'etale the diagonal
$U \to U \times_X U$ is an open immersion, see
Morphisms of Spaces, Lemmas
\href{https://stacks.math.columbia.edu/tag/06CR}{lemma etale unramified (Tag 06CR)} and
\href{https://stacks.math.columbia.edu/tag/05W1}{lemma diagonal unramified morphism (Tag 05W1)}.
Hence $W_p$ is an open and closed subspace of
$U^{p + 1} = U \times_X \ldots \times_X U$. The action of $S_{p + 1}$
on $W_p$ is free as we've thrown out the fixed points of the action.
By
Lemma \ref{lemma-product-is-tensor-product}
we see that
$$
(f_!\underline{\mathbf{Z}})^{\otimes p + 1} =
f^{p + 1}_!\underline{\mathbf{Z}} = (W_p \to X)_!\underline{\mathbf{Z}}
\oplus Rest
$$
where $f^{p + 1} : U^{p + 1} \to X$ is the structure morphism.
Looking at stalks over a geometric point $\overline{x}$ of $X$
we see that
$$
\left(
\bigoplus\nolimits_{\overline{u} \mapsto \overline{x}} \mathbf{Z}
\right)^{\otimes p + 1}
\longrightarrow
(W_p \to X)_!\underline{\mathbf{Z}}_{\overline{x}}
$$
is the quotient whose kernel is generated by all tensors
$1_{\overline{u}_0} \otimes \ldots \otimes 1_{\overline{u}_p}$
where $\overline{u}_i = \overline{u}_j$ for some $i \not = j$.
Thus the quotient map
$$
(f_!\underline{\mathbf{Z}})^{\otimes p + 1}
\longrightarrow
\wedge^{p + 1}f_!\underline{\mathbf{Z}}
$$
factors through $(W_p \to X)_!\underline{\mathbf{Z}}$, i.e., we get
$$
(f_!\underline{\mathbf{Z}})^{\otimes p + 1}
\longrightarrow
(W_p \to X)_!\underline{\mathbf{Z}}
\longrightarrow
\wedge^{p + 1}f_!\underline{\mathbf{Z}}
$$
This already proves that $\Hom(K^p, \mathcal{F})$ is (functorially) a
subgroup of
$$
\Hom((W_p \to X)_!\underline{\mathbf{Z}}, \mathcal{F}) = \mathcal{F}(W_p)
$$
To identify it with the $S_{p + 1}$-anti-invariants we have to prove that
the surjection $(W_p \to X)_!\underline{\mathbf{Z}}
\to \wedge^{p + 1}f_!\underline{\mathbf{Z}}$ is the maximal
$S_{p + 1}$-anti-invariant quotient. In other words, we have to show that
$\wedge^{p + 1}f_!\underline{\mathbf{Z}}$ is the quotient of
$(W_p \to X)_!\underline{\mathbf{Z}}$ by the subsheaf generated by
the local sections $s - \text{sign}(\sigma)\sigma(s)$ where $s$ is
a local section of $(W_p \to X)_!\underline{\mathbf{Z}}$.
This can be checked on the stalks, where it is clear.
\end{proof}

\begin{lemma}
\label{lemma-twist}
Let $S$ be a scheme. Let $W$ be an algebraic space over $S$.
Let $G$ be a finite group acting freely on $W$.
Let $U = W/G$, see
Properties of Spaces, Lemma \href{https://stacks.math.columbia.edu/tag/071S}{lemma quotient (Tag 071S)}.
Let $\chi : G \to \{+1, -1\}$ be a character.
Then there exists a rank 1 locally free sheaf of $\mathbf{Z}$-modules
$\underline{\mathbf{Z}}(\chi)$ on $U_\etale$ such that for every
abelian sheaf $\mathcal{F}$ on $U_\etale$ we have
$$
H^0(W, \mathcal{F}|_W)^\chi =
H^0(U, \mathcal{F} \otimes_{\mathbf{Z}} \underline{\mathbf{Z}}(\chi))
$$
\end{lemma}

\begin{proof}
The quotient morphism $q : W \to U$ is a $G$-torsor, i.e., there exists
a surjective \'etale morphism $U' \to U$ such that
$W \times_U U' = \coprod_{g \in G} U'$ as spaces with $G$-action over $U'$.
(Namely, $U' = W$ works.) Hence $q_*\underline{\mathbf{Z}}$ is a finite
locally free $\mathbf{Z}$-module with an action of $G$. For any
geometric point $\overline{u}$ of $U$, then we get $G$-equivariant
isomorphisms
$$
(q_*\underline{\mathbf{Z}})_{\overline{u}}
= \bigoplus\nolimits_{\overline{w} \mapsto \overline{u}} \mathbf{Z}
= \bigoplus\nolimits_{g \in G} \mathbf{Z} = \mathbf{Z}[G]
$$
where the second $=$ uses a geometric point
$\overline{w}_0$ lying over $\overline{u}$ and
maps the summand corresponding to $g \in G$ to the summand
corresponding to $g(\overline{w}_0)$. We have
$$
H^0(W, \mathcal{F}|_W) =
H^0(U, \mathcal{F} \otimes_\mathbf{Z} q_*\underline{\mathbf{Z}})
$$
because
$q_*\mathcal{F}|_W = \mathcal{F} \otimes_\mathbf{Z} q_*\underline{\mathbf{Z}}$
as one can check by restricting to $U'$. Let
$$
\underline{\mathbf{Z}}(\chi) =
(q_*\underline{\mathbf{Z}})^\chi \subset
q_*\underline{\mathbf{Z}}
$$
be the subsheaf of sections that transform according to $\chi$. For
any geometric point $\overline{u}$ of $U$ we have
$$
\underline{\mathbf{Z}}(\chi)_{\overline{u}} =
\mathbf{Z} \cdot \sum\nolimits_g \chi(g) g
\subset
\mathbf{Z}[G] = (q_*\underline{\mathbf{Z}})_{\overline{u}}
$$
It follows that $\underline{\mathbf{Z}}(\chi)$ is locally free of
rank 1 (more precisely, this should be checked after restricting to $U'$).
Note that for any $\mathbf{Z}$-module $M$ the $\chi$-semi-invariants
of $M[G]$ are the elements of the form $m \cdot \sum\nolimits_g \chi(g) g$.
Thus we see that for any abelian sheaf $\mathcal{F}$ on $U$ we have
$$
\left(\mathcal{F} \otimes_\mathbf{Z} q_*\underline{\mathbf{Z}}\right)^\chi
=
\mathcal{F} \otimes_\mathbf{Z} \underline{\mathbf{Z}}(\chi)
$$
because we have equality at all stalks. The result of the lemma follows by
taking global sections.
\end{proof}


\begin{lemma}
\label{lemma-alternating-spectral-sequence}
Let $S$ be a scheme. Let $f : U \to X$ be a surjective, \'etale, and
separated morphism of algebraic spaces over $S$. For $p \geq 0$ set
$$
W_p = U \times_X \ldots \times_X U \setminus \text{all diagonals}
$$
(with $p + 1$ factors) as in Lemma \ref{lemma-compute}.
Let $\chi_p : S_{p + 1} \to \{+1, -1\}$ be the sign character.
Let $U_p = W_p/S_{p + 1}$ and $\underline{\mathbf{Z}}(\chi_p)$ be as in
Lemma \ref{lemma-twist}.
Then the spectral sequence of
Lemma \ref{lemma-alternating-cech-to-cohomology}
has $E_1$-page
$$
E_1^{p, q} =
H^q(U_p, \mathcal{F}|_{U_p} \otimes_\mathbf{Z} \underline{\mathbf{Z}}(\chi_p))
$$
and converges to $H^{p + q}(X, \mathcal{F})$.
\end{lemma}

\begin{proof}
Note that since the action of $S_{p + 1}$ on $W_p$ is over $X$ we do
obtain a morphism $U_p \to X$. Since $W_p \to X$ is \'etale and since
$W_p \to U_p$ is surjective \'etale, it follows
that also $U_p \to X$ is \'etale, see
Morphisms of Spaces, Lemma \href{https://stacks.math.columbia.edu/tag/03XT}{lemma etale local (Tag 03XT)}.
Therefore an injective object of
$\textit{Ab}(X_\etale)$ restricts to an injective object of
$\textit{Ab}(U_{p, \etale})$, see
Cohomology on Sites, Lemma \href{https://stacks.math.columbia.edu/tag/03F3}{lemma cohomology of open (Tag 03F3)}.
Moreover, the functor
$\mathcal{G} \mapsto
\mathcal{G} \otimes_\mathbf{Z} \underline{\mathbf{Z}}(\chi_p))$
is an auto-equivalence of $\textit{Ab}(U_p)$, whence transforms injective
objects into injective objects and is exact (because
$\underline{\mathbf{Z}}(\chi_p)$ is an invertible
$\underline{\mathbf{Z}}$-module). Thus given an injective resolution
$\mathcal{F} \to \mathcal{I}^\bullet$ in $\textit{Ab}(X_\etale)$
the complex
$$
\Gamma(U_p,
\mathcal{I}^0|_{U_p} \otimes_\mathbf{Z} \underline{\mathbf{Z}}(\chi_p))
\to
\Gamma(U_p,
\mathcal{I}^1|_{U_p} \otimes_\mathbf{Z} \underline{\mathbf{Z}}(\chi_p))
\to
\Gamma(U_p,
\mathcal{I}^2|_{U_p} \otimes_\mathbf{Z} \underline{\mathbf{Z}}(\chi_p))
\to \ldots
$$
computes
$H^*(U_p,
\mathcal{F}|_{U_p} \otimes_\mathbf{Z} \underline{\mathbf{Z}}(\chi_p))$.
On the other hand, by
Lemma \ref{lemma-twist}
it is equal to the complex of $S_{p + 1}$-anti-invariants in
$$
\Gamma(W_p, \mathcal{I}^0) \to
\Gamma(W_p, \mathcal{I}^1) \to
\Gamma(W_p, \mathcal{I}^2) \to \ldots
$$
which by
Lemma \ref{lemma-compute}
is equal to the complex
$$
\Hom(K^p, \mathcal{I}^0) \to
\Hom(K^p, \mathcal{I}^1) \to
\Hom(K^p, \mathcal{I}^2) \to \ldots
$$
which computes
$\Ext^*_{\textit{Ab}(X_\etale)}(K^p, \mathcal{F})$.
Putting everything together we win.
\end{proof}
\end{document}
